%------------------------------------------------------------------------------%
% Luis A. D'Afonseca
%
%------------------------------------------------------------------------------%
\documentclass[a4paper,12pt,fleqn]{article}

\usepackage{style_avaliacao}

\ShowAnswers

\newcommand{\D}[2]{\frac{\partial {#1}}{\partial {#2}}}
\newcommand{\DS}[2]{\frac{\partial^2 {#1}}{\partial {#2}^2}}
\newcommand{\DM}[3]{\frac{\partial^2 {#1}}{\partial {#2} \partial {#3}}}


\newcommand{\vetor}[2]{\left(\begin{array}{c} #1 \\ #2 \end{array}\right)}

\newcommand{\veci}{\bm{i}}
\newcommand{\vecj}{\bm{j}}
\newcommand{\veck}{\bm{k}}

%------------------------------------------------------------------------------%
\begin{document}

\makeHeader{CFVVI}{Exame Especial -- Turma 6}{18/09/2024}

% 01
%------------------------------------------------------------------------------%
\exercise{25}
Calcule o limite solicitado, ou prove que o limite não existe,
\(
  \lim_{(x, y)\to(0, 0)} \frac{x-y}{x+y}
\)

\begin{answer}
  Analisando a função percebemos que se fizermos $y=x$ o limite restrito
  a curva será
  \[
    L_1
    = \lim_{(x, y)\to(0, 0)} \left(\frac{x-y}{x+y} \evalat{y=x}{}\right)
    = \lim_{x\to 0} \frac{0}{2x} = 0
  \]
  Por outro lado, escolhendo $y=2x$ temos
  \[
    L_2
    = \lim_{(x, y)\to(0, 0)} \left(\frac{x-y}{x+y} \evalat{y=2x}{}\right)
    = \lim_{x\to 0} \frac{-x}{3x}
    = \frac{-1}{3}
  \]
  Como $L_1\neq L_2$ o limite não existe.
  \clearpage
\end{answer}

% 02
%------------------------------------------------------------------------------%
\exercise{25}
Encontre linearização da função
\(
f(x, y) = \sqrt{y-x}
\)
no ponto \((1, 2)\).

\begin{answer}
  A linearização de $f$ no ponto \((1, 2)\) é a função
  \begin{align*}
    L(x, y)
    & = f(x_0, y_0) + f_x(x_0, y_0)(x-x_0) + f_y(x_0, y_0)(y-y_0) \\
    & = f(1, 2) + f_x(1, 2)(x-1) + f_y(1, 2)(y-2)
  \end{align*}
  Precisamos das derivadas parciais de $f$
  \[
    f_x(x, y)
    = \D{f}{x}
    = \D{}{x}\left[\sqrt{y-x}\,\right]
    = \D{}{x}\left[(y-x)^{\sfrac{1}{2}}\right]
    = \frac{1}{2}(y-x)^{\sfrac{-1}{2}}\D{}{x}(y-x)
    = \frac{-1}{2\sqrt{y-x}}
  \]
  \[
    f_y(x, y)
    = \D{f}{y}
    = \D{}{y}\left[\sqrt{y-x}\,\right]
    = \D{}{y}\left[(y-x)^{\sfrac{1}{2}}\right]
    = \frac{1}{2}(y-x)^{\sfrac{-1}{2}}\D{}{y}(y-x)
    = \frac{1}{2\sqrt{y-x}}
  \]
  Avaliando $f$ e suas derivadas no ponto $(1, 2)$
  \[
    f(1, 2)
    = \sqrt{2-1}
    = \sqrt{1}
    = 1
  \]
  \[
    f_x(1, 2)
    = \frac{-1}{2\sqrt{2-1}}
    = \frac{-1}{2}
  \]
  \[
    f_y(1, 2)
    = \frac{1}{2\sqrt{2-1}}
    = \frac{1}{2}
  \]
  Assim
  \begin{align*}
    L(x, y)
    & = f(1, 2) + f_x(1, 2)(x-1) + f_y(1, 2)(y-2) \\
    & = 1 -\frac{1}{2}(x-1) + \frac{1}{2}(y-2) \\
    & = 1 -\frac{x}{2} + \frac{1}{2} + \frac{y}{2} - 1 \\
    & = -\frac{x}{2} + \frac{y}{2} + \frac{1}{2}
  \end{align*}
  \clearpage
\end{answer}

% 03 ex 20 pg 270 res pg 838
%------------------------------------------------------------------------------%
\exercise{25}
Encontre e classifique os pontos críticos da função
\(
  f(x, y) = x^4 + y^4 + 4xy
\)

\begin{answer}
  Precisamos das derivadas parciais de $f$
  \[
    \D{f}{x}
    = \D{}{x}\left[x^4 + y^4 + 4xy\right]
    = 4x^3 + 4y
  \]
  \[
    \D{f}{y}
    = \D{}{y}\left[x^4 + y^4 + 4xy\right]
    = 4y^3 + 4x
  \]
  então
  \[
    \nabla f = \vetor{4x^3 + 4y}{4y^3 + 4x}
  \]

  A função é um polinômio, então possui derivadas em todos os pontos do plano.
  Assim os pontos críticos são apenas os pontos onde as derivadas parciais são zero,
  \(\nabla f = 0\)
  \[
    4x^3 + 4y = 0
    \qquad\text{e}\qquad
    4y^3 + 4x = 0
  \]
  ou, simplificando,
  \[
    x^3 + y = 0
    \qquad\text{e}\qquad
    y^3 + x = 0
  \]
  Isolando $y$ na primeira equação, \(y = -x^3\), e substituindo na segunda, temos
  \begin{align*}
    y^3 + x                 & = 0 \\
    \left(-x^3\right)^3 + x & = 0 \\
    -x^9 + x                & = 0 \\
    x^9 - x                 & = 0 \\
    x(x^8 - 1)              & = 0
  \end{align*}
  As soluções dessa equação são $x=0$, $x=1$ ou $x=-1$.
  Se $x=0$ temos $y=0$,
  se $x=1$ temos $y=-1$ e
  se $x=-1$ temos $y=1$.
  Portanto, os pontos críticos são
  \[
    (x_1, y_1) = (0, 0),
    \qquad\qquad
    (x_2, y_2) = (1, -1)
    \qquad\text{e}\qquad
    (x_3, y_3) = (-1, 1)
  \]

  Precisamos das derivadas parciais de segunda ordem de $f$
  \[
    \DS{f}{x}
    = \D{}{x}\left[4x^3 + 4y\right]
    = 12x^2
  \]
  \[
    \D{f}{y}
    = \D{}{y}\left[4y^3 + 4x\right]
    = 12y^2
  \]
  \[
    \DM{f}{x}{y}
    = \D{}{x}\D{f}{y}
    = \D{}{x}\left[4y^3 + 4x\right]
    = 4
  \]
  então
  \[
    H = \left(\begin{array}{cc}
      12x^2 & 4 \\
      4 & 12y^2
    \end{array}\right)
  \]


  Para classificar os pontos críticos precisamos avaliar o discriminante nos
  pontos críticos.

  Considerando o ponto \((x_1, y_1) = (0, 0)\)
  \[
    f_{xx}(0, 0) = 0
  \]
  \[
    f_{yy}(0, 0) = 0
  \]
  \[
    f_{xy}(0, 0) = 4
  \]
  \[
    D_1
    = f_{xx}(0, 0)f_{yy}(0, 0) - f^2_{xy}(0, 0)
    = 0 \times 0 - 4^2
    = -16 < 0
  \]
  Portanto, o ponto $(0, 0)$ é um ponto de sela.

  Considerando o ponto \((x_2, y_2) = (1, -1)\)
  \[
    f_{xx}(1, -1) = 12
  \]
  \[
    f_{yy}(1, -1) = 12
  \]
  \[
    f_{xy}(1, -1) = 4
  \]
  \[
    D_2
    = f_{xx}(1, -1)f_{yy}(1, -1) - f^2_{xy}(1, -1)
    = 12 \times 12 - 4^2
    = 144 -16
    = 128 > 0
  \]
  Portanto, o ponto $(1, -1)$ é um máximo ou mínimo local. Como \(f_{xx}(1, -1) = 12 > 0\)
  o ponto é um ponto de mínimo local.

  Considerando o ponto \((x_3, y_3) = (-1, 1)\)
  \[
    f_{xx}(-1, 1) = 12
  \]
  \[
    f_{yy}(-1, 1) = 12
  \]
  \[
    f_{xy}(-1, 1) = 4
  \]
  \[
    D_3
    = f_{xx}(-1, 1)f_{yy}(-1, 1) - f^2_{xy}(-1, 1)
    = 12 \times 12 - 4^2
    = 144 -16
    = 128 > 0
  \]
  Portanto, o ponto $(-1, 1)$ é um máximo ou mínimo local. Como \(f_{xx}(-1, 1) = 12 > 0\)
  o ponto é um ponto de mínimo local.
  \clearpage
\end{answer}

% 04
%------------------------------------------------------------------------------%
\exercise{25}
Encontre as dimensões da caixa retangular fechada com máximo volume
que pode ser inscrita na esfera unitária.

\begin{answer}
  Queremos maximizar a função
  \[
    f(x, y, z) = (2x)(2y)(2z) = 8xyz
  \]
  sujeita à
  \[
    g(x, y, z) = x^2 + y^2 + z^2 = 1
  \]
  Sabemos também que $x$, $y$ e $z$ precisam sere positivos para que o volume $f$
  seja positivo.

  Para aplicarmos multiplicadores de Lagrange precisamos das derivadas
  parciais de $f$ e $g$
  \[
    f_x(x, y, z) = 8yz \qquad\qquad
    f_y(x, y, z) = 8xz \qquad\qquad
    f_z(x, y, z) = 8xy
  \]
  \[
    g_x(x, y, z) = 2x \qquad\qquad
    g_y(x, y, z) = 2y \qquad\qquad
    g_z(x, y, z) = 2z
  \]
  Vamos agora resolver o sistema \(\nabla f = \lambda\nabla g\) e \(g=1\)
  \[
  \begin{cases}
    8yz = 2\lambda x \\
    8xz = 2\lambda y \\
    8xy = 2\lambda z \\
    x^2 + y^2 + z^2 = 1
  \end{cases}
  \]
  Simplificando
  \[
  \begin{cases}
    4yz = \lambda x \\
    4xz = \lambda y \\
    4xy = \lambda z \\
    x^2 + y^2 + z^2 = 1
  \end{cases}
  \]

  Isolando o $\lambda$ na primeira equação
  \[
    \lambda = \frac{4yz}{x}
  \]
  e substituindo na segunda
  \[
    4xz
    = \lambda y
    = \frac{4yz}{x} y
    = \frac{4y^2z}{x}
    \qquad\Rightarrow\qquad
    x^2 = y^2
  \]
  Substituindo na terceira
  \[
    4xy = \lambda z
    = \frac{4yz}{x}z
    = \frac{4yz^2}{x}
    \qquad\Rightarrow\qquad
    x^2 = z^2
  \]
  Assim a terceira equação se torna
  \begin{align*}
    x^2 + y^2 + z^2 & = 1                      \\[2mm]
    3x^2            & = 1                      \\[2mm]
    x^2             & = \frac{1}{3}            \\[2mm]
    x               & = \frac{1}{\sqrt{3}}
  \end{align*}
  Portanto
  \[
    x = y = z = \frac{1}{\sqrt{3}}
  \]
  \clearpage
\end{answer}

%------------------------------------------------------------------------------%
\end{document}
%------------------------------------------------------------------------------%
